\scanpage{292}
et que ces suites et la suite constante de valeur 1 forment une base de \unsure{l'espace} vectoriel $E_r$ de toutes les suites complexes de période $r$. Montrer que la série
\[
u(s)_n c_n
\]
est alternée à valeurs absolues décroissantes tendant vers zéro, donc convergente. En déduire une autre démonstration de b).

d) Soit $z$ un nombre complexe, et posons $u_n=z^n$ pour $n\quad1$. Montrer que la suite des sommes partielles $S_n$ de (1) reste bornée si et seulement si on a
\[
(5)\qquad z\quad1\ ,\quad z\ne1
\]
(auquel cas la suite des $S_n$ est majorée par $2\,\frac{z}{1-z}$).
\deleted{\unsure{Montrer que si la série $c_n$ est divergente, alors}}
\inserted{En déduire que} la condition (5) implique que la série \inserted{$c_nz^n$}\deleted{$\unsure{s_n|u_n}$} est convergente. Montrer qu'elle est même absolument convergente pour $z\quad1$, tandis que pour un $z$ tel que $z=1$, $z\ne1$, on a convergence absolue si et seulement si la \deleted{suite} série $c_n$ est convergente. Enfin,
\deleted{Application à la \unsure{convergence} \illegible{3} \unsure{Montrer que}}
\inserted{$A$ désignant une partie}\deleted{\unsure{pour toute partie $A$}} de l'ensemble \inserted{$E$} des nombres complexes satisfaisant (5), pour que la convergence de la série $c_nz^n$ soit uniforme pour $z\quad A$, il suffit qu'il existe un nombre $a\quad0$ tel que $z\quad A$ implique $1-z\quad a$. Montrer que cette condition est aussi nécessaire si la série $c_n$ est divergente (tandis que si $c_n$ est convergente, la convergence est uniforme pour l'ensemble des nombres complexes $z$ tels que $z\quad1$). \deleted{\unsure{E}}

e) Montrer que la fonction $f(z)$ sur l'ensemble $E$ des nombres complexes $z$ satisfaisant (5), représentée par la série
\[
(6)\qquad f(z)=c_nz^n
\]
est continue. Elle admet une limite finie pour $z=1$ si et seulement si la \deleted{suite} série $c_n$ est convergente, auquel cas la fonction $f(z)$ définie par la formule (6) est définie et continue \deleted{pour} sur tout le disque unité $z\quad1$.

f) Appliquer les résultats de d) et e) à la série trigonométrique
\[
c_ne^{ix}\quad,
\]
et aux séries
\[
c_n\cos nx\ ,\quad c_n\sin nx\quad;
\]
en particulier, ces séries sont convergentes pour tout $x$ réel n'appartenant pas à l'ensemble \deleted{$\unsure{\mathbb R}$} $2\quad\underline{\mathrm Z}$ (i.e. $x$ pas de la forme $2\ n$, avec $n$ entier), et représentent des fonctions continues de $x$ pour $x\quad\underline{\mathrm R}-2\quad\underline{\mathrm Z}$.

g) Montrer que la suite $(z^n)_{n\quad1}$ ($z$ étant un nombre complexe) est périodique de période $r$ si et seulement si $z$ est une racine $r$.ème de l'unité. \deleted{\unsure{En déduire}} Déduire de ceci et et b) que lorsque $z$ est une racine de l'unité la série \deleted{$\unsure{u_n}$} $c\ z$ est convergente, sauf dans le cas où $z=1$ et $c_n$ \unsure{est}

\margin{293}
